% =============================================================================
% Capítulo 1 · Introducción
% Fuente: tesis/_capitulos/01_introduccion.qmd (aprobado en la iteración 2). Se traslada
% sin afirmaciones nuevas; las cifras de los datos salen de cifras.tex.
% =============================================================================
\chapter{Introducción}
\label{cap:introduccion}

\section{Planteamiento y descripción del problema}
\label{sec:planteamiento}

El ciclo financiero es la fluctuación de mediano plazo entre fases de expansión y contracción
del crédito, de los precios de activos y de la disposición a asumir riesgos.
\textcite{drehmann_borio_tsatsaronis_2012} lo distinguen del ciclo económico por su duración, de
8 a 30 años, y documentan que sus picos coinciden con las crisis bancarias sistémicas. No se
observa de forma directa: se manifiesta en el movimiento conjunto de muchas variables, cuya
estructura de relaciones puede cambiar en el tiempo. Medirlo es, por tanto, un problema de
reducción dimensional sobre series de tiempo multivariadas.

Para el Perú, el antecedente más directo es
\textcite{camiz_cerda_atoche_2026}\verificar{manuscrito fechado el 25 de abril de 2026, sin
datos de revista, volumen ni DOI; confirmar con el Dr.~Camiz su estado de publicación.}, que
continúa la tesis de maestría de Atoche Murrieta \pendiente{Atoche Murrieta (2022), tesis de
maestría, PUCP; sin ficha}. Ese trabajo aplica el Análisis de Componentes Principales (PCA), la
Clasificación Factorial Jerárquica (HFC) y el Análisis Evolutivo de Componentes Principales
(EPCA) a 17 indicadores mensuales del Banco Central de Reserva del Perú (BCRP). Deja abiertas,
sin embargo, cuatro cuestiones: no evalúa el aprendizaje profundo, no clasifica fases del ciclo
contrastadas con episodios fechados de forma independiente, su muestra termina en 2020 y fija
el ancho de la ventana del EPCA sin un criterio objetivo.

% Fuente de las cifras: cifras.tex (herramientas/generar_cifras.py), a partir de las
% dimensiones que escribe notebooks/05_pca.ipynb en data/results/.
Esta tesis analiza \nSeries{} series mensuales del BCRP de crédito, liquidez, tasas, precios,
tipo de cambio, bolsa, riesgo soberano, actividad y reservas. El análisis oficial, sobre las
series sin transformar (decisión D1), cubre \nMesesNiveles{} meses, de \fechaInicioNiveles{} a
\fechaFinDatos{}. La vía transformada tiene \nMesesTransformadas{} meses, de
\fechaInicioTransformadas{} a \fechaFinDatos{}, porque la diferenciación pierde el primer mes.
Ambas vías se ejecutan con el mismo código (decisión D2), y la forma de combinarlas está
pendiente de discusión con el asesor.

\textbf{Problema de investigación:} ¿cómo identificar, fechar y caracterizar los regímenes del
ciclo financiero peruano entre \anioInicioDatos{} y \anioFinDatos{} a partir de la estructura
multivariada de los indicadores del BCRP, integrando los métodos exploratorios evolutivos de la
escuela de Camiz con autoencoders y redes LSTM, sin renunciar a la interpretabilidad ni al
principio de que la estructura emerge de los datos?

\section{Justificación de la investigación}
\label{sec:justificacion}

\textbf{Metodológica.} La tesis amplía la cadena PCA--HFC--EPCA de tres maneras. Fecha los
regímenes con la partición óptima de
\textcite{fisher_1958}\verificar{páginas (789--798), solo en fuentes secundarias; número y
DOI ya verificados. Ficha preliminar, texto no leído: solo se cita su aporte general.}, como hacen
\textcite{camiz_roig_2011}\verificar{el nombre del editor de las actas (O.~Papú) se infirió por
analogía; el año figura como 2010/2011 en la ficha.}. Propone un criterio objetivo para el ancho
de la ventana del EPCA. E incorpora autoencoders y redes LSTM anidados en la cadena clásica,
porque un autoencoder lineal con pérdida cuadrática genera el subespacio de los primeros
componentes principales
\parencite{bourlard_kamp_1988,baldi_hornik_1989}\verificar{Bourlard y Kamp (1988), número del
fascículo; ficha preliminar, texto no consultado. Baldi y Hornik (1989), solo el DOI; el texto
completo ya se leyó.}; véase la \cref{prop:ae-lineal}. Como los antecedentes de aprendizaje profundo
usan conjuntos de datos mucho mayores
\parencite{hinton_salakhutdinov_2006,tuhin_et_al_2025}\verificar{Hinton y Salakhutdinov (2006),
se procesó el material suplementario; Tuhin et al. (2025), lista completa de autores.},
documentar cómo se adaptan a series macrofinancieras cortas es en sí un resultado.

\textbf{Económica y social.} \textcite{drehmann_borio_tsatsaronis_2012} acuñan la noción de
«recesión inconclusa»: si la política ignora la expansión del crédito y de los precios de los
inmuebles, el ajuste se posterga y la crisis posterior resulta más severa. Un indicador de
régimen transparente y reproducible podría servir como insumo exploratorio para los análisis de
estabilidad financiera del BCRP y de la SBS.

\section{Objetivos}
\label{sec:objetivos}

\subsection{Objetivo general}

Formalizar y evaluar una metodología exploratoria evolutiva, que integre el Análisis de
Componentes Principales, la Clasificación Factorial Jerárquica y el Análisis Evolutivo de
Componentes Principales con autoencoders y redes LSTM, para detectar, fechar y caracterizar los
regímenes del ciclo financiero peruano entre \anioInicioDatos{} y \anioFinDatos{} a partir de
indicadores macrofinancieros del BCRP.

\subsection{Objetivos específicos (compromisos firmes)}
\label{sec:oe}

\begin{enumerate}[label=\textbf{OE\arabic*.}, ref=OE\arabic*, leftmargin=*]
  \item\label{oe:estado-arte} Sistematizar el estado del arte sobre el ciclo financiero y sobre
    la reducción dimensional de series de tiempo multivariadas.
  \item\label{oe:datos} Construir y versionar el conjunto de series mensuales del BCRP, con
    trazabilidad de códigos, transformaciones y exclusiones.
  \item\label{oe:clasicos} Aplicar PCA, HFC y EPCA en las dos vías del diseño y contrastar los
    resultados con los de \textcite{camiz_cerda_atoche_2026}.
  \item\label{oe:ventana} Formalizar un criterio objetivo para el ancho de ventana del EPCA,
    basado en la similitud entre ventanas consecutivas (coeficiente RV o distancia de
    Procrustes).
  \item\label{oe:robustez-hfc} Evaluar la robustez de los grupos de la HFC mediante remuestreo
    por bloques y correlación cofenética.
  \item\label{oe:dimension} Triangular el número de componentes con la vara quebrada, el
    criterio de Rencher y el análisis paralelo de Horn
    \parencite{horn_1965}\verificar{ficha preliminar: se leyó el resumen, no el texto; solo se
    cita el aporte general.} \pendiente{fuentes de la vara quebrada y del criterio de Rencher}.
  \item\label{oe:dl} Diseñar autoencoders y redes LSTM adaptados a la escala de los datos, con
    pesos atados, preentrenamiento por capas, regularización y validación por bloques temporales
    y \emph{walk-forward}.
  \item\label{oe:incertidumbre} Cuantificar la incertidumbre del error de reconstrucción
    mediante Monte Carlo Dropout
    \parencite{gal_ghahramani_2016}\verificar{editores del volumen; PDF no consultado.}.
  \item\label{oe:datacion} Fechar regímenes con la partición óptima de Fisher sobre
    trayectorias de representaciones (primer valor propio del EPCA, \emph{scores} del PCA o
    códigos latentes) y clasificarlos en expansión, contracción o transición.
  \item\label{oe:episodios} Contrastar regímenes y fases con episodios fechados por fuentes
    externas, nunca por las propias series ni por un indicador derivado de ellas.
  \item\label{oe:simulado} Comparar los métodos sobre un contexto simulado, con regímenes
    conocidos por construcción, y calibrar con ello el número de factores y de clases, el ancho
    de ventana y los pesos de los análisis evolutivos.
  \item\label{oe:comparacion} Comparar los métodos clásicos y los de aprendizaje profundo según
    interpretabilidad, estabilidad temporal, capacidad descriptiva y robustez.
\end{enumerate}

\subsection{Objetivos exploratorios condicionales}

Se abordarán si los resultados de los objetivos firmes lo justifican:

\begin{enumerate}[label=\textbf{OC\arabic*.}, ref=OC\arabic*, leftmargin=*]
  \item\label{oc:ehfc} Evaluar la viabilidad de una HFC evolutiva (EHFC) que aplique la
    jerarquía completa en cada ventana móvil, en lugar de la aplicación secuencial de HFC y EPCA
    de la literatura.
  \item\label{oc:statis} Explorar el STATIS dual como detector independiente de quiebres
    estructurales, siguiendo a
    \textcite{camiz_maulucci_roig_2010}\verificar{datos de edición de las actas (ISBN o DOI).}.
  \item\label{oc:anticipar} Explorar la capacidad del LSTM para anticipar transiciones, con
    validación estrictamente fuera de muestra y declarando su escaso poder estadístico.
  \item\label{oc:extension} Extender la muestra a años anteriores a \anioInicioDatos{}, si la
    homogeneidad de las series lo permite, para evaluar la crisis financiera global fuera de
    muestra.
\end{enumerate}

\section{Hipótesis}
\label{sec:hipotesis}

El estudio es exploratorio y no formula hipótesis contrastables en sentido estricto, por tres
razones. Los métodos centrales describen la muestra sin postular un modelo que la genere, de
modo que no hay una hipótesis nula. Los pocos episodios de referencia no dan poder estadístico.
Y los resultados no se generalizan a otras economías. Se formulan, en cambio, \textbf{hipótesis
de trabajo}, que se evalúan por la convergencia de la evidencia en las dos vías. Cada una indica
qué resultado contaría en su contra.

\begin{description}[style=nextline, leftmargin=1.2em]
  \item[HT1 (no estacionariedad estructural).] La estructura de correlación no es
    constante entre \anioInicioDatos{} y \anioFinDatos{}: hay intervalos estables separados por
    transiciones. \emph{En contra:} que las ventanas sucesivas resulten tan similares como
    cabría esperar bajo una estructura constante.
  \item[HT2 (organización en grupos).] Las series se organizan en pocos grupos
    homogéneos que se corresponden con los bloques identificados por
    \textcite{camiz_cerda_atoche_2026}: salud económica y financiera, tasas de interés y riesgo
    del mercado financiero. Como el conjunto de series no es idéntico, la correspondencia se
    evalúa por bloque. \emph{En contra:} grupos inestables bajo remuestreo o sin
    correspondencia con esos bloques.
  \item[HT3 (correspondencia con episodios).] Los regímenes se corresponden,
    dentro de una tolerancia declarada de antemano, con episodios fechados externamente.
    \emph{En contra:} cortes sin relación con ningún episodio, o que dependan de la regla que
    fija su número.
  \item[HT4 (complementariedad).] Los autoencoders y las redes LSTM aportan
    información sobre la dinámica y la atipicidad del sistema que el PCA no capta por
    construcción. \emph{En contra:} que no mejoren a su anclaje lineal en ningún criterio fuera
    de muestra, o que sus resultados cambien con la semilla aleatoria.
\end{description}
